\documentclass[12pt]{article}

%%%% tyhle 3 radky jsou pro dvoustranku %%%%%%%%%%%%%%%%

%\documentclass[10pt,twocolumn]{article}  ěščřžýáíé
%\usepackage[landscape, left=2cm, right=2cm]{geometry}
%\columnsep 6pc  %   Space between columns

\usepackage[margin=3cm]{geometry}
%\setlength{\oddsidemargin}{0mm}
%\setlength{\evensidemargin}{0mm}
%\addtolength{\textheight}{55mm}
%\addtolength{\textwidth}{43mm}
%\setlength{\topmargin}{0mm}
%\setlength{\headsep}{4mm}
%\voffset=-20mm
%\hoffset=-20mm
%\columnseprule=.01pt


%\usepackage[iso8859-2]{inputenc}
\usepackage[utf8]{inputenc}
%\usepackage[babel]{czech}

\usepackage{amssymb}
\usepackage{enumitem}   % pro enumerate s pismenky
\usepackage{hyperref}   % use for hypertext links, including those to external documents and URLs


\setlength\parindent{0pt}



\pagestyle{empty}

\begin{document}

{\Large{Determinanty  verze 5}} 
-- uvažujeme zatím pouze \textbf{reálné} matice.


{\section{Homomorfismy $R^* \rightarrow R^*$ }} 

\textbf{Definice:} $R^* = R - \{0\}$, $R^*$ uažujeme jako multiplikativní grupu
\vspace{1mm}

Homomorfismus $R^* \rightarrow R^*$ je každá reálná funkce $g$, pro kterou platí $g(x\,y) = g(x)\,g(y)$. 
\vspace{4mm}

\textbf{Úloha:} Chceme zjistit, jakou formu mají všechny homomorfismy $R^* \rightarrow R^*$.
\vspace{2mm}

\textbf{Věta 1:} Nechť $g$ je homomorfismus $R^* \rightarrow R^*$. \\
Pak platí, že buď $g(x) = |x|^\alpha$, nebo $g(x) = sign(x)\,|x|^\alpha$, kde $\alpha \in R$.

\vspace{2mm}
Důkaz: viz tvůj text, kapitoly 2  a 3 - tohle by se mělo nějak sloučit do jednoho důkazu pro $R^*$ jako multiplikativní grupu, bez odkazu na ten aditivní zápis v kap. 2, aby se vůbec nezaváděl.

$\Box$
\vspace{5mm}


%%%%%%%%%%%%%%%%%%%%%%%%%%%%%%%%%%%%%%%%%%%%%%%%%%%%
{\section{Restrikce homomorfismů $GL(n) \rightarrow R^*$ na $GD(n)$}} 

\textbf{Definice:} 
\vspace{1mm}
$GL(n)$ je grupa všech reálných regulárních matic.
\vspace{1mm}
Homomorfismus $GL(n) \rightarrow R^*$ je každá funkce $f$, pro kterou platí \begin{equation}
f(AB) = f(A)\,f(B) .
\end{equation}

\textbf{Poznámka:} pro homomorfismy platí 
\begin{align*}
f(P\,A\,P^{-1}) &= f(P)\,f(A)\,f(P^{-1}) = f(P)\,f(P^{-1})\,f(A)=f(P\,P^{-1})\,f(A) \\
  &=  f(I)\,f(A) = f(A),
\end{align*}
kde jsme využili faktu, že $f(I)=1$ (viz Appendix).

\vspace{2mm}

Grupu všech regulárních diagonálních matic označíme jako $GD(n)$,\\
 je to podgrupa $GL(n)$.
\vspace{4mm}

\textbf{Úloha:} Chceme zjistit, jak vypadá restrikce libovolného homomorfismu\\ $GL(n) \rightarrow R^*$ na podgrupu $GD(n)$. 
\vspace{2mm}

\textbf{Věta 2:} Nechť $F: GL(n) \rightarrow R^*$ je homomorfismus a
 $f: GD(n) \rightarrow R^*$ \\
 je jeho restrikce na $GD(n)$. Pak $f$ se dá vyjádřit jako homomorfismus \\
 $g: R^* \rightarrow R^*$ tak, že $f(D) \equiv g(\Pi\, d_i)$.
\vspace{2mm}

Důkaz: 
\vspace{1mm}
Použijeme permutační matice $P_i$, které prohodí 1. a $i$-tý řádek. $P_i$ vznikne z $I$ prohozením 1. a $i$-tého řádku, resp. prohozením 1. a $i$-tého sloupce (vyjde totéž). Platí $P_i^{-1} = P_i$, násobení zleva touto maticí prohodí řádky, násobení zprava prohodí sloupce; pro diagonální matici $D$ operace $P_i D P_i^{-1}$ prohodí 1. a $i$-tý prvek na diagonále.

\begin{align*}
f(D) &=  f(diag(d_1, 1, \dots 1))\, f(diag(1, d_2, 1, \dots 1)) \dots \, f(diag(1, \dots 1,  d_n)) \\
  &= f(diag(d_1, 1, \dots 1))\,f(P_2\,diag(d_2, 1, \dots 1)\,P_2^{-1})\dots \\
  &= \dots f(P_n\,diag(d_n, 1, \dots 1)\,P_n^{-1}) \\
  &= f(diag(d_1, 1, \dots 1))\,f(diag(d_2, 1, \dots 1))\dots f(diag(d_n, 1, \dots 1))\, \\
  &= f(diag(\Pi\, d_i, 1, \dots 1)) \equiv g(\Pi\, d_i).
\end{align*}

$g$ je homomorfismus: $g(x)$ je def. jako $g(x) = f(diag(x, 1, \dots 1)0$, takže
\begin{align*}
g(x\,y) &= f(diag(xy, 1, \dots 1)) 
    = f(diag(x, 1, \dots 1)\, diag(y, 1, \dots 1)) \\
  &= f(diag(x, 1, \dots 1))\,f(diag(y, 1, \dots 1)) \\
  &= g(x)\,g(y)
\end{align*}
$\Box$

\vspace{3mm}

Jestliže $f$ je homomorfismus $GD(n) \rightarrow R^*$, pak dle Věty 2
se dá vyjádřit jako \\
$f(D) \equiv g(\Pi\, d_i)$, kde $g$ je homomorfismus $R^* \rightarrow R^*$. Odtud dle Věty 1 plyne, že  \\
buď 
\begin{equation}
\label{h1}
f(D) \equiv g(\Pi\, d_i) = |\Pi\, d_i|^\alpha \, ,
\end{equation}
nebo 
\begin{equation}
\label{h2}
f(D) \equiv g(\Pi\, d_i)= sign(\Pi\, d_i)\,|\Pi\, d_i|^\alpha \, ,
\end{equation}
kde $\alpha \in R$.
\vspace{6mm}


%%%%%%%%%%%%%%%%%%%%%%%%%%%%%%%%%%%%%%%%%%%%%%%%%%%%
{\section{Homomorfismy $GL(n) \rightarrow R^*$ }} 


\textbf{Úloha:} Hledáme všechny homomorfismy $GL(n) \rightarrow R^*$.

\vspace{3mm}

\textbf{Věta 3:} Každý homomorfismus na $GL(n)$ je jednoznačně určen svou restrikcí\\
 na podgrupu $GD(n)$.

\vspace{2mm}
Důkaz: 
\vspace{1mm}
\begin{enumerate}[label=(\alph*)]
\item 
If $C$ is diagonizable, then $C=P\,D\,P^{-1}$, 
where $D$ is diagonal and $P$ is invertible, and
\begin{align*}
f(C) &= f(P\,D\,P^{-1}) = f(P)\,f(D) \,f(P^{-1}) = f(P)\,f(P^{-1})\,f(D)\\
 &=  f(P\,P^{-1})\,f(D) = f(D).
\end{align*}
Diagonální matice $D$ je určena jednoznačně až na pořadí prvků na diagonále (jsou na ní vlastní čísla matice $C$), tj. číslo $\Pi\, d_i$ je jednoznačné.

\item 
For general $A$, use polar decomposition $A=H\,U$, where both $H$ and $P$ are diagonizable. Then using the result from a) above:
\begin{equation}
\label{AD}
f(A) = f(H\,U) = f(H)\,f(U) = f(D_H)\,f(D_U) = f(D_H\,D_U) = f(D).
\end{equation}
Zde již prvky matice  $D$ nejsou určeny jednoznačně, protože $D_H$ a $D_U$ nemají určené pořadí prvků na diagonále, ale číslo $\Pi\, d_i$ už je jednoznačné.
\end{enumerate}
Matice $D$ tedy není určena jednoznačně, avšak číslo $\Pi\, d_i$ je. Ze vztahů \ref{h1} a \ref{h2} tedy plyne, že rozšíření z $GD(n)$ na $GL(n)$ je jednoznačné.

$\Box$
\vspace{3mm}

Závěr: všechny homomorfismy $GL(n) \rightarrow R^*$ mají tvar 

\vspace{1mm}
buď 
\begin{equation}
\label{h11}
f(A) = f(D) = |\Pi\,  d_i|^\alpha \, ,
\end{equation}
nebo 
\begin{equation}
\label{h22}
f(A) = f(D) = sign(\Pi\,  d_i)\,|\Pi\,  d_i|^\alpha \, ,
\end{equation}
kde $\alpha \in R$ a $D$ je matice z \ref{AD} (přesně řečeno, kterákoliv z matic získaných tímto postupem). 
\vspace{2mm}



%\newpage
%%%%%%%%%%%%%%%%%%%%%%%%%%%%%%%%%%%%%%%%%%%%%%%%%%%%
{\section{Homogenní homomorfismy $GL(n) \rightarrow R^*$ }} 

\textbf{Definice:} 
\vspace{1mm}
Homomorfismus $f: GL(n) \rightarrow R^*$ se nazývá homogenní (stupně $n$), jestliže\\
pro všechna $\lambda \in R^*$ platí 
\begin{equation}
\label{homo}
f(\lambda I) = \lambda^n .
\end{equation}
\textbf{Poznámka:} 
\vspace{1mm}
Z této definice plyne $f(\lambda A) = f((\lambda I) A) =f(\lambda I) \,f(A) = \lambda^n f(A)$.
\vspace{3mm}

\textbf{Definice:} 
\vspace{1mm}
Homomorfismus $f: GL(n) \rightarrow R^*$ se nazývá pozitivně homogenní, jestliže rovnost \ref{homo} \\
platí pro všechna $\lambda > 0$.
\vspace{1mm}
\textbf{Poznámka:} 
\vspace{1mm}
Pro $n$ sudé obě definice splývají.\\
Pro $n$ liché je pozitivní homogenita pouze nutnou podmínkou pro homogenitu.
\vspace{2mm}


\vspace{3mm}
\textbf{Věta 4.1:} Oba homomorfismy \ref{h11} a \ref{h22} jsou pozitivně homogenní na $GL(n)$,\\
 právě když $\alpha = 1$. 
\vspace{2mm}
Důkaz -- ověření rovnosti \ref{homo} pro $\lambda > 0$:
\begin{equation}
\label{hi1}
f(\lambda I) = |\Pi\, d_i|^\alpha = |\lambda^n|^\alpha  = (\lambda^n)^\alpha\, ,
\end{equation}
resp.
\begin{equation}
\label{hi2}
f(\lambda I) = sign(\Pi\, d_i)\,|\Pi\, d_i|^\alpha 
= sign(\lambda^n)\,|\lambda^n|^\alpha = (\lambda^n)^\alpha  \, ,
\end{equation}
takže $f$ je pozitivně homogenní, právě když $\alpha = 1$.  

$\Box$

\vspace{3mm}
\textbf{Věta 4.2:} Každý homogenní homomorfismus $f: GL(n) \rightarrow R^*$ má tvar\\  pro sudé $n$ buď
\begin{equation}
\label{h111}
f(A) = f(D) = |\Pi\, d_i|\, ,
\end{equation}
nebo
\begin{equation}
\label{h222}
f(A) = f(D) = \Pi\, d_i\, ,
\end{equation}
kde $D$ je libovolná matice z \ref{AD} . 
Pro liché $n$ má tvar pouze \ref{h222} .

\vspace{2mm}
Důkaz: 
\vspace{1mm}
Pro sudé $n$ plyne z Věty 4.1 a z toho, že pro sudé $n$ z pozitivní homogenity plyne homogenita. 
Pro liché $n$ může nastat případ $\lambda^n < 0$, který vylučuje variantu \ref{h111}. 

$\Box$
\vspace{4mm}


%%%%%%%%%%%%%%%%%%%%%%%%%%%%%%%%%%%%%%%%%%%%%%%%%%%%
{\section{Definice determinantu na $GL(n)$ }} 

\vspace{4mm}

Na $GD(n)$ definujme determinant jako homogenní homomorfismus $GD(n) \rightarrow R^*$.\\
 Z Věty 4.2 plyne, že pro liché $n$ je definice jednoznačná a že 
\begin{equation}
\label{detD}
\det(D) = \Pi\, d_i \, .
\end{equation}
Tato definice souhlasí se standardní definicí determinantu.
\vspace{1mm}
Dodatečná podmínka pro sudé $n$, aby definice byla jednoznačná ?  Nevím.

Mohli bychom například přidat podmínku, že $f(P_2) < 0$, kde $P_2$ je matice\\
 z důkazu Věty 2. Nebo podmínku $f(diag(-1, 1, \dots 1)) = -1$.


\vspace{5mm}
Na $GL(n)$ definujme determinant jako rozšíření z grupy $GD(n)$, které je \\
dle Věty 3 jednoznačné a je dáno předpisem
\begin{equation}
\label{detA}
\det(A) = \det(D) \, ,
\end{equation}
kde $D$ je libovolná matice z \ref{AD} . 
\vspace{4mm}


\textbf{Věta 5:} Výše uvedená definice determinantu na $GL(n)$ souhlasí se standardní definicí.

\vspace{2mm}
Důkaz: 
\vspace{1mm}
Determinant dle standardní definice a dle naší definice se shodují na $GD(n)$.
\vspace{1mm}
Determinant na $GD(n)$ dle standardní definice i dle naší definice představují rozšíření z $GD(n)$ na $GL(n)$.
\vspace{1mm}
Věta 3 říká, že rozšíření homomorfismu z $GD(n)$ na $GL(n)$ je jednoznačné. \\
Takže se obě definice musí shodovat.

$\Box$
\vspace{6mm}




A teď se to ještě musí projít a ověřit, že $R^*$ se dá nahradit $C^*$, a pak ještě z nějaké spojitosti odvodit, že to platí i pro neregulární matice.


\vspace{3mm}

%%%%%%%%%%%%%%%%%%%%%%%%%%%%%%%%%%%%%%%%%%%%%%%%%%%%
{\section{Appendix} 

\vspace{3mm}

Platí, že složení dvou homomorfismů je homomorfismus, a že obraz jednotkového prvku je jednotkový prvek - viz např.
\vspace{1mm}
{\url{<https://sites.math.washington.edu/~greenber/GroupHomomorphisms.pdf>}}
\vspace{2mm}


\end{document}
